\chapter{Patrones transversales}
\label{sec:patrones}

Hay cinco patrones que se repiten en todo el código del proyecto.
Cada uno resuelve un problema concreto; reconocerlos acelera la
lectura de cualquier módulo nuevo.

\section{Typestate}

Codificar invariantes en el sistema de tipos para que estados
inválidos no compilen. Se usa prominentemente en
\filepath{core/matcher-ffi}:

\begin{lstlisting}[language=Rust]
pub struct GmGraph<S: State> {
    ptr: *mut c_void,
    _state: PhantomData<S>,
}

pub struct Building;
pub struct Finalized;

impl GmGraph<Building> {
    pub fn add_edge(&mut self, ...) -> Result<()> { ... }
    pub fn finalize(self) -> GmGraph<Finalized> { ... }
}

impl GmGraph<Finalized> {
    pub fn search(&self, ...) -> GmResults { ... }
}
\end{lstlisting}

\code{search} sólo existe en el estado \code{Finalized}; intentar
buscar en un \code{GmGraph<Building>} es un error de compilación,
no de runtime.

\section{RAII sobre handles C}

Todo puntero crudo que retorna \filepath{libgraphmatch} queda
envuelto en una struct Rust que implementa \code{Drop} llamando a
la función \code{gm_*\_free} correspondiente. Esto garantiza:

\begin{itemize}
    \item Liberación determinística en el scope donde se crea el
    handle.
    \item Liberación correcta ante \code{panic} (drop se ejecuta
    durante el unwind).
    \item Imposibilidad de ``double free'' — el tipo consume su
    contenido al liberar.
\end{itemize}

Aplica a \code{GmGraph}, \code{GmMatcher} y \code{GmResults}. Antes
del rediseño, \code{GmMatcher} liberaba manualmente dentro de una
función, lo que abría una ventana unsafe ante panic.

\section{Canonical API}

Una pieza de lógica se implementa \textbf{una sola vez},
en \filepath{platform/api}, y todos los transportes (HTTP, IPC,
stdin, MCP) son adaptadores sobre esa misma lógica. Lo opuesto a
la estructura anterior, donde Tauri, CLI y API reimplementaban cada
uno partes del mismo flujo.

La regla operativa: si te encuentras copiando código de un crate
de aplicación a otro, ese código debería estar viviendo en
\filepath{platform/api}.

\section{Model Gate}

Separar la adquisición de un modelo (archivo ONNX) de su uso
mediante una capa de validación entre ambos:

\begin{lstlisting}[language=Rust]
pub struct ModelGate { ... }

impl ModelGate {
    pub fn validate(path: &Path, expected: &ModelSpec)
        -> Result<ValidatedModel, GateError> { ... }
}

pub struct ValidatedModel { /* inference-ready */ }
\end{lstlisting}

Sólo \code{ValidatedModel} puede pasarse al inferencer. Si la
validación falla (checksum, dimensiones, versión), el caller
decide qué hacer: el patrón no fuerza un crash, pero sí impide
pasar un modelo no validado accidentalmente.

Uso actual: \filepath{core/gnn/common/model_gate.rs} decide si el
scorer GNN se habilita.

\section{Registry de extensiones}

Donde la plataforma necesita extensibilidad dinámica (sumar un
parser, un entity type del DSL, una tool MCP) se usa un registry:

\begin{lstlisting}[language=Rust]
pub trait LogParser: Send + Sync {
    fn name(&self) -> &str;
    fn can_parse(&self, sample: &[u8]) -> bool;
    fn parse(&self, bytes: &[u8]) -> Result<Vec<Event>>;
}

pub struct ParserRegistry {
    parsers: Vec<Box<dyn LogParser>>,
}

impl ParserRegistry {
    pub fn register(&mut self, p: Box<dyn LogParser>) { ... }
    pub fn select(&self, sample: &[u8]) -> Option<&dyn LogParser>
      { ... }
}
\end{lstlisting}

Variantes del mismo patrón viven en \filepath{platform/dsl}
(registry de extensiones de gramática), \filepath{platform/mcp}
(registry de tools) y \filepath{platform/local-llm} (registry de
backends).

El contrato es siempre: trait estable, registry concreto, discovery
en tiempo de arranque. \emph{No} hay autodescubrimiento por
reflexión; todos los registros son explícitos en código.

\section{Por qué estos cinco y no otros}

Son los patrones que \emph{cruzan capas}. Hay muchas otras
convenciones locales (por ejemplo, el estilo de errores en
\filepath{platform/api} usando \code{thiserror}) pero no afectan
cómo se diseñan módulos nuevos en otras capas.

Estos cinco sí: cada uno tiene un costo (aumenta el código
boilerplate o exige conocimiento previo) compensado por un
beneficio estructural (imposibilita bugs, garantiza unicidad,
desacopla extensibilidad). Cuando aparece un módulo nuevo que
necesita invariantes en tipos, manejo de recursos externos, lógica
pública, validación de artefactos o extensibilidad dinámica, se
alcanza primero a uno de estos patrones antes de inventar otro.
